\documentclass[12pt]{article}
\usepackage[utf8]{inputenc}
\usepackage[english]{babel}
\usepackage[a4paper,left=2cm,right=2cm,top=2cm,bottom=2cm]{geometry}
\usepackage{indentfirst}
\usepackage{appendix}
\usepackage{libertine}
\usepackage{titlesec}
\usepackage{eso-pic}
\usepackage{fancyhdr}
\usepackage{graphicx}
\usepackage{wrapfig}
\usepackage{float}
\usepackage{caption}
\usepackage{subcaption}
\usepackage{color}
\usepackage{setspace}
\usepackage{multicol}
\usepackage[T1]{fontenc}
\usepackage[hyphens]{url}
\usepackage[backend=bibtex,style=alphabetic]{biblatex}
\addbibresource{bibliography.bib}
\usepackage{bm}
\usepackage{enumitem}
\usepackage{siunitx}
\usepackage{textcomp}
\usepackage{amsmath}
\usepackage{amssymb}
\usepackage{amsfonts}
\usepackage{amsthm}
\usepackage{commath}
\usepackage{cancel}
\usepackage{pdflscape}
\usepackage{svg}
\usepackage[linesnumbered,ruled,vlined]{algorithm2e}
\usepackage{comment}
\usepackage{placeins}
\usepackage{listings}
\usepackage[linkbordercolor=white]{hyperref}
\usepackage{glossaries}
\makeglossaries
\loadglsentries{glossaries.tex}

\lstset{literate=%
    {à}{{\`a}}1 {â}{{\^a}}1 {é}{{\'e}}1 {è}{{\`e}}1 {ê}{{\^e}}1 {ë}{{\"e}}1 {î}{{\^i}}1 {ï}{{\"i}}1 {ô}{{\^o}}1 {ö}{{\"o}}1 {ù}{{\`u}}1 {û}{{\^u}}1 {ü}{{\"u}}1 {ç}{{\c{c}}}1 {Ç}{{\c{C}}}1{―}{{\textemdash}}1
    {’}{{'}}1 {‘}{{`}}1 {“}{{``}}1 {”}{{''}}1
}
\urlstyle{same}
\pagestyle{fancy}
\fancyhead[L]{\leftmark}
\fancyhead[R]{}
\renewcommand{\subsectionmark}[1]{}

% Configuration de \subsubsubsection
\titleclass{\subsubsubsection}{straight}[\subsection]
\newcounter{subsubsubsection}[subsubsection]
\renewcommand\thesubsubsubsection{\thesubsubsection.\arabic{subsubsubsection}}
\renewcommand\theparagraph{\thesubsubsubsection.\arabic{paragraph}}

\titleformat{\subsubsubsection}
  {\normalfont\normalsize\bfseries}{\thesubsubsubsection}{1em}{}
\titlespacing*{\subsubsubsection}
{0pt}{3.25ex plus 1ex minus .2ex}{1.5ex plus .2ex}

\makeatletter
\renewcommand\paragraph{\@startsection{paragraph}{5}{\z@}%
  {3.25ex \@plus1ex \@minus.2ex}%
  {-1em}%
  {\normalfont\normalsize\bfseries}}
\renewcommand\subparagraph{\@startsection{subparagraph}{6}{\parindent}%
  {3.25ex \@plus1ex \@minus .2ex}%
  {-1em}%
  {\normalfont\normalsize\bfseries}}
\def\toclevel@subsubsubsection{4}
\def\toclevel@paragraph{5}
\def\toclevel@subparagraph{6} % Correction du doublon ici
\def\l@subsubsubsection{\@dottedtocline{4}{7em}{4em}}
\def\l@paragraph{\@dottedtocline{5}{10em}{5em}}
\def\l@subparagraph{\@dottedtocline{6}{14em}{6em}}
\makeatother

\setcounter{secnumdepth}{4}
\setcounter{tocdepth}{4}

\setlength{\headheight}{14.5pt}
\setlength{\parskip}{1em}
\renewcommand{\baselinestretch}{0.6}
\renewcommand{\headrulewidth}{1pt}
\newcommand{\HRule}{\rule{\linewidth}{0.3mm}}
\newcommand\BackgroundPic{%
        \put(0,0){%
                \parbox[b][\paperheight]{\paperwidth}{%
                        \vfill
                        \centering
                        \vfill}}}

% -----------------------------------------------------

\begin{document}
\AddToShipoutPicture*{\BackgroundPic}
\begin{titlepage}
  \begin{sffamily}
\noindent
\begin{minipage}{0.32\textwidth}
    \raggedright
    \includegraphics[height=2.0cm]{faculty_logo.png}
\end{minipage}%
\begin{minipage}{0.32\textwidth}
    \centering
    \includegraphics[height=2.0cm]{science.png}
\end{minipage}%
\begin{minipage}{0.32\textwidth}
    \raggedleft
    \includegraphics[height=2.0cm]{bioscience.png}
\end{minipage}
\vspace*{5cm}

  \begin{center}

    % Title
    \textsc{\huge C003698A - Design Project}\\[1cm]

    \HRule \\[0.7cm]

    \textsc{\Huge Exploring the value of omics data for individual cancer
prognosis using multimodal AI} \\[0.4cm]

    \HRule \\[1.2cm]

    % Author and supervisor
    \begin{minipage}{0.5\textwidth}
    \begin{flushleft} \large
    \emph{Author:}\\
    Romain \textsc{Liefferinckx} - 02611641\\
    Dane Marc Yap \textsc{Bagaoisan} - \\
    Lotte \textsc{Bonte} - 02305915\\
    Victoria \textsc{Likhacheva} - 01906610\\
    \end{flushleft}
    \end{minipage}
    \begin{minipage}{0.4\textwidth}
    \begin{flushright} \large
    \emph{Supervisors:} \\
    Fiebe \textsc{Denijs}\\
    Willem \textsc{Waegaman}\\
    \end{flushright}
    \end{minipage}

    \vfill

    \textsc{ UGent – Dept. Data Analysis and Mathematical Modelling}

    % Faculty logo above academic year
    \includegraphics[height=4.0cm]{logo_UGent.png}\\[1.0cm]

    % Bottom of the page
    {\large Academic year 2026-2027}

  \end{center}

  \end{sffamily}
\end{titlepage}

\clearpage

\tableofcontents
\newpage

% -----------------------------------------------------
\nocite{*}
\begin{multicols}{2}
\section*{Abstract}
\addcontentsline{toc}{section}{Abstract}
Predicting cancer progression and survival can support physicians in making decisions about patient
care and treatment. Multimodal AI can improve these predictions by combining clinical information,
imaging and molecular profiles. However, collecting additional data can be costly, and their value may
differ between patients. A particular omics assay may improve predictions for one patient while
providing redundant information for another. The goal of this project is to investigate whether
uncertainty estimates can help determine when additional omics data are needed for an individual
patient and which type would be most informative.
\section{Introduction}

\section{Data}

\section{Uncertainty Estimation}

\section{Methods for deep learning with multimodal data}

\subsection{Early Fusion (Data Level Fusion)}
Unified Multimodal Transformers process heterogeneous inputs within a single, 
end-to-end architecture by projecting all modalities such as image patches, clinical narratives, 
and omics profiles into a shared sequence of \glspl{token_embedding}. 
While naive token concatenation often incurs a prohibitive quadratic computational complexity 
($\mathcal{O}(N^2)$) relative to the combined sequence length, 
architectures utilizing latent bottlenecks (such as Perceiver or Perceiver IO) 
mitigate this challenge. They employ cross attention mechanisms to map high dimensional input sequences into a fixed size latent space, 
where deep self attention blocks perform unified cross modal reasoning. 
Missing modalities or missing clinical features are natively handled via dynamic attention masking: tokens corresponding to unobserved data are assigned a large negative bias 
($-\infty$ pre softmax), preventing unobserved modalities from altering the latent representations while preserving a constant parameter layout.

\subsection{Intermediate Fusion (Feature Level Fusion)}
Intermediate Fusion bridges early and late approaches by processing each modality through domain dedicated feature extractors 
(encoders) to yield high level latent representations. 
These representations are subsequently merged within the intermediate layers of the network using dedicated interaction mechanisms, 
such as cross attention or tensor fusion. 
This paradigm strikes a balance by preserving domain specific inductive biases while enabling deep, 
nonlinear interactions across modalities. 
Missing data is typically managed by substituting empty feature embeddings combined with binary presence indicator masks, 
allowing the intermediate fusion modules to dynamically adapt to incomplete observations.

\subsection{Late Fusion (Decision Level Fusion)}
Late Fusion operates by processing each modality through an independent, 
domain specific expert model up to the point of final inference. 
Predictions from these individual models 
(e.g., class probabilities or risk scores) 
are subsequently aggregated at the decision level using strategies such as weighted averaging, 
voting, or meta classification. 
Because features are never directly mixed, 
this paradigm accommodates missing data naturally: if a modality is unavailable for a given sample, its corresponding expert pipeline is simply omitted from the final aggregation. 
However, the principal limitation of Late Fusion lies in its inability to model cross modal feature interactions during representation learning, 
restricting its capacity to capture complex inter modality dependencies.

\subsection{Multimodal Mixture of Experts}
The Multimodal Mixture of Experts 
(MoE) architecture utilizes a dynamic routing mechanism to direct incoming data streams toward specialized subnetworks 
("experts") conditioned on the input modalities present. Instead of activating all parameters for every sample, 
the gating network evaluates the available data and selectively assigns computation to the most relevant domain experts. 
In scenarios involving incomplete observations, 
the router bypasses nonviable execution paths, ensuring that missing modalities do not trigger inactive expert networks. 
This conditional computation paradigm achieves sublinear scaling with respect to model parameters, 
balancing domain specific precision with flexible multisource integration.

\section{Conclusions}


\newpage
\glsaddall
\printglossaries

\newpage
\printbibliography
\newpage
\section*{Appendix}
\appendix
\end{multicols}
\end{document}
